\section{Coverage of the entire tetrahedron}\label{sec:coverage}

The shared-boundary identity gives an angular sector in each family for
every point of $K$. The remaining task is to prove that at least one of the
two selected triangles has sufficient radial extent.

\begin{proposition}\label{prop:coverage}
For every $0<\eps\le\eta/2$, the $2n$ cylinders
\eqref{eq:first-family}--\eqref{eq:second-family}, with parameters
\eqref{eq:parameters}--\eqref{eq:cutoffs}, cover the closed tetrahedron $K$.
\end{proposition}
\begin{proof}
\emph{Choosing angular sectors.} Fix $(x,y,Ht)\in K$ and write $s=1-t$.
The intercept along $v_j$ is $(0,y_0,Ht_0)$, where
\begin{equation}\label{eq:lower-intercept}
 t_0=t-\eps x\beta_j,\qquad y_0=y-\eps x\alpha_j.
\end{equation}
Consider the finite sequence
\[
 F_i=q_it+\eps x\frac{1-q_i^2}{4}\qquad(0\le i\le n).
\]
Its endpoints are $F_0=-t\le y\le t=F_n$.
Choose the first index $i\ge1$ for which $F_i\ge y$, and set $j=i-1$.
Such an index exists because $F_n\ge y$. If $i=1$, then $F_j=F_0\le y$;
if $i>1$, minimality gives $F_j<y$. Thus
$F_j\le y\le F_{j+1}$, including when $y$ equals an extreme endpoint.
Using \eqref{eq:shared-boundaries}, these inequalities are exactly
\begin{equation}\label{eq:lower-sector}
 q_jt_0\le y_0\le q_{j+1}t_0.
\end{equation}
Subtracting their outer terms gives $0\le(q_{j+1}-q_j)t_0=\Delta t_0$.
Since $\Delta>0$, they force $t_0\ge0$.
Hence the selected first-family cylinder covers the point unless $t_0>T_j$.
This argument does not assume that the sequence $F_i$ is monotone.

For the second family, use the sequence
\[
 G_i=q_is-\eps y\frac{1-q_i^2}{4},\qquad G_0=-s\le x\le s=G_n.
\]
The same first-crossing choice gives an index $k$ for which
\begin{equation}\label{eq:upper-intercept}
 s_0=s+\eps y\beta_k\ge0,\qquad x_0=x+\eps y\alpha_k,\qquad
 q_ks_0\le x_0\le q_{k+1}s_0.
\end{equation}
This cylinder covers unless $s_0>T_k$.

\emph{Comparing the radial cutoffs.} It remains to show that the two selected radial cutoffs cannot both fail.
Suppose both selected cylinders miss the point.
Then $t_0,s_0>1/2$, and we can set
\[
 q=y_0/t_0\in[q_j,q_{j+1}],\qquad
 p=x_0/s_0\in[q_k,q_{k+1}].
\]
Since $|x|+|y|\le1$, \eqref{eq:coefficient-bounds} gives
\[
 t_0+s_0=1-\eps x\beta_j+\eps y\beta_k\le1+\eps/2.
\]
Thus each of $t_0,s_0$ lies between $1/2$ and $1/2+\eps/2$.
Using $y=t_0q+\eps x\alpha_j$ and $x=s_0p-\eps y\alpha_k$, we obtain
\begin{equation}\label{eq:rough-angles}
 |y-q/2|\le\eps,\qquad |x-p/2|\le\eps.
\end{equation}

Since $|\beta_j-q/2|\le\Delta/4$ and
$|\beta_k-p/2|\le\Delta/4$, \eqref{eq:rough-angles} gives
$|\beta_j-y|,|\beta_k-x|\le\eps+\Delta/4$.
Thus the intercept heights $t_0=t-\eps x\beta_j$ and
$s_0=s+\eps y\beta_k$ have opposite first-order changes
$-\eps xy$ and $+\eps xy$. This cancellation is the reason for choosing
$\beta_j$ close to $q/2$.

To control the remaining terms exactly, write $a=t_0-1/2$ and
$b=s_0-1/2$. The two radial excesses satisfy
\begin{align}
 a(1-\eps xq)+b(1+\eps yp)
 &=\eps^2(x^2\alpha_j+y^2\alpha_k)\notag\\
 &\quad-\eps x(\beta_j-q/2)+\eps y(\beta_k-p/2).
 \label{eq:radial-budget}
\end{align}
For completeness, $a+b=-\eps x\beta_j+\eps y\beta_k$, while the
intercept equations give
$q a=y-\eps x\alpha_j-q/2$ and
$p b=x+\eps y\alpha_k-p/2$. Substitute these three identities into the
left side of \eqref{eq:radial-budget}; the two mixed terms cancel.

For $0<\eps<1$, both factors multiplying $a$ and $b$ are at least
$1-\eps>0$. Since $a>\eps^2M_j$, $b>\eps^2M_k$ and $|x|+|y|\le1$, dividing
\eqref{eq:radial-budget} by $\eps^2$ yields
\begin{equation}\label{eq:failure-inequality}
 0>(M_j+M_k)(1-\eps)-\frac{\Delta}{4\eps}
                   -(x^2\alpha_j+y^2\alpha_k).
\end{equation}

We bound the last term without losing uniformity at the boundary.
From \eqref{eq:rough-angles} and $|x|,|y|\le1$, $|p|,|q|\le1$,
\[
 |x^2-p^2/4|,\ |y^2-q^2/4|\le\tfrac32\eps.
\]
For a point inside its interval,
\[
 \left|\alpha_j-\frac{1+q^2}{4}\right|\le\Delta/2,\qquad
 \left|\alpha_k-\frac{1+p^2}{4}\right|\le\Delta/2.
\]
Consequently
\begin{align*}
 x^2\alpha_j+y^2\alpha_k
 &\le\frac{p^2+q^2+2p^2q^2}{16}+\frac32\eps+\frac\Delta4\\
 &\le d(p)+d(q)+\frac32\eps+\frac\Delta4.
\end{align*}
The second inequality explains our choice of radial enlargement:
\[
 d(p)+d(q)-\frac{p^2+q^2+2p^2q^2}{16}
 =\frac{(p^2-q^2)^2}{16}\ge0.
\]
The cutoff definitions give
\[
 d(p)+d(q)+2\eta\le M_j+M_k\le\frac14+2\eta.
\]
Using $\Delta\le\eps^2$, the right side of \eqref{eq:failure-inequality} is therefore at least
\begin{equation}\label{eq:uniform-coverage-gap}
 2\eta-(2+2\eta)\eps-\frac{\eps^2}{4}.
\end{equation}
For $0<\eps\le\eta/2$ this is at least
$\eta-17\eta^2/16>0$, a contradiction.
All sector and cutoff inequalities are non-strict membership conditions. The argument applies to every point of the closed tetrahedron, including its faces, edges and vertices.
\end{proof}

We retain the numerical estimate with a variable margin for the later
construction in Section~\ref{sec:angular-cubic}.

\begin{corollary}[Quantitative coverage criterion]
\label{cor:quantitative-coverage}
Let $0<\tau<1$ and $\eta\ge0$. Choose an integer $n$ with
$\Delta=2/n\le\tau^2$, put $q_j=-1+j\Delta$, and define
$\alpha_j,\beta_j$ by \eqref{eq:axis-coefficients}. Use the triangles and
directions \eqref{eq:first-family}--\eqref{eq:second-family}, with
$\eps$ replaced by $\tau$ and with cutoffs
\[
 T_j=\frac12+\tau^2\left(\eta+\max_{q_j\le q\le q_{j+1}}d(q)\right).
\]
These $2n$ cylinders cover $K$ whenever
\begin{equation}\label{eq:variable-margin-budget}
 2\eta-(2+2\eta)\tau-\frac{\tau^2}{4}\ge0.
\end{equation}
\end{corollary}
\begin{proof}
The angular first-crossing argument is unchanged. If both selected caps
failed, their intercept heights would exceed $1/2$ and have sum at most
$1+\tau/2$. Thus the rough-angle estimates and the radial budget remain
valid with the same constants, independently of $\eta$. The only later
use of the margin is through
$d(p)+d(q)+2\eta\le M_j+M_k\le1/4+2\eta$.
Consequently the same calculation as
\eqref{eq:failure-inequality}--\eqref{eq:uniform-coverage-gap} gives
$0>2\eta-(2+2\eta)\tau-\tau^2/4$, contradicting
\eqref{eq:variable-margin-budget}.
\end{proof}

For the fixed parameters of the main construction, the coverage threshold
$\eps\le1/2000$ will also guarantee strict area saving. We now bound
the area remainder explicitly, so no further choice of a smaller tilt
is needed.
